\documentclass[../main.tex]{subfiles}
\begin{document}
\section{序論}
\label{sec:introduction}
% 本文をここに記述する。
\subsection{(物理目標)}
どんな問いがあるか
\subsection{J-PARC チャームバリオン分光実験}
E50が提案されていて準備が進んでいる、どんな実験で、どんなことを目的としているか
\subsection{実験ビームライン}
高運動量ビームラインについて
\subsection{MARQスペクトロメータとそのデータ収集システム}
MARQの現状と連続読み出し式DAQについて述べ、なぜ連続読み出し式DAQ(+ オンラインフィルタ on it)なのか述べる
\begin{tcolorbox}[
  title=執筆メモ,
  colback=yellow!8,
  colframe=black!30,
  breakable
]
\begin{itemize}
  \item MARQスペクトロメータの現状
  \begin{itemize}
    \item TOFシンチレータ
    \item RPC
    \item DC
    \item FT
    \item ThAC
    \item B-RICH
    \item RICH
  \end{itemize}
  \item 合計チャンネル数(あとで確認)
  \item DAQ
  \begin{itemize}
    \item フルストリーミングDAQ
    \item その利点(トリガービジーによるデッドタイムがないシステムにできること、物理の面ではアンバイアスな測定ができること、実験遂行の面ではハードウェアトリガーを作るためのコストがかからないこと)
    \item 留意点(ビームスルーイベントの割合、賢くデータを貯めるにはオンラインでのデータ削減が有効であること)
    \item 用いるシステム
    \begin{itemize}
      \item NestDAQ
      \item nestdaq-user-impl
      \item 実験個別のHLF
    \end{itemize}
  \end{itemize}
\end{itemize}
\end{tcolorbox}

\subsection{研究目的}
\begin{tcolorbox}[
  title=執筆メモ,
  colback=yellow!8,
  colframe=black!30,
  breakable
]
\begin{itemize}
  \item E50を遂行するためにどれくらいの計算資源が必要かを見積もる
  \item ソフトウェアによるデータリダクション
  \item どんなオンラインフィルタが有効か?
  \item それにはどれくらいの計算資源が必要か?
\end{itemize}
\end{tcolorbox}


連続読み出し式DAQシステムのテスト実験T103のデータを用いてスケールすることで見積もる
\end{document}
